% !TeX root = master.tex
% !TeX spellcheck = de_DE

\chapter{Aufbereiten der Daten}

    Was waren die Ursprungsdaten, und wie habe ich diese Aufgebereitet (Subsampling und so)

\section{Ursprungsdaten}
    Rant darüber wie schlecht das Ursprungspaper Dokumentiert war



\subsection{Literatur}
    Alle wissenschaftliche Literatur, die man verwendet, wird im Literatur-Verzeichnis aufgeführt.
    Dieses wird automatisch von Bibtex erzeugt.
    Dazu muss zunächst ein Eintrag in der Bibliographie-Datenbank im bib-Ordner angelegt werden.
    Man findet z.\,B. mit Google Scho\-lar für wissenschaftliche Veröffentlichungen fertige Zitationsangaben im Bibtex-Format.

    Wenn ich eine Veröffentlichung zitiere, sieht das so aus.~\cite{Krauthoff2017a}
    Man kann auch, insbesondere bei längeren Quellen wie Büchern, eine Seitenzahl ergänzen.~\cite[S.~5]{Krauthoff2017a}
    Es ist auch möglich, mehrere Quellen auf einmal zu nennen.\cite{Krauthoff2017a,dun:j:argument-acceptability}


\section{CleaningPipeline}

    Was habe ich gemacht um an meine Batches/Subsamples zu kommen?

